\section{The Camera-Mounted IMU as a Pose Sensor}
\label{sec:imu}

\subsection{One-Second Rigidity Test}
Adjacent 10-Hz planar-PnP rotations are too noisy for a rigidity decision, so
the proposed test compares relative rotations over 0.8--1.2~s. With
$\boldsymbol\phi_I$ the integrated IMU rotation vector and
$\boldsymbol\phi_C=\operatorname{Log}^{\vee}
[{\bf R}_{BC}^{\top}(t_a){\bf R}_{BC}(t_b)]$,
alternating increments estimate the constant extrinsic rotation by
\begin{equation}
 \hat{\bf R}_{CI}=\arg\min_{{\bf R}\in SO(3)}
 \sum_k\|{\bf R}\boldsymbol\phi_{I,k}-\boldsymbol\phi_{C,k}\|^2 ,
\end{equation}
where $\operatorname{Log}^{\vee}$ is the rotation-vector logarithm of
$SO(3)$. The rotation is fitted on one half of the increments, and the
remaining increments measure per-axis correlation, magnitude
correlation, scale and RMS residual. The test requires deliberate pitch, yaw
and roll excitation with at least 1$^\circ$ standard deviation per axis, and
acceptance requires each held-out axis correlation $\geq0.75$, magnitude
correlation $\geq0.90$ and scale in $[0.85,1.15]$; correlation $\geq0.90$ on
all axes is marked high quality. On board-visible intervals of 5--30~s the
same comparison yields an empirical gyroscope drift rate, reported per axis
because planar-PnP sensitivity is anisotropic.

\subsection{Label-Leakage Test}
To quantify leakage, a second label set replaces the PnP rotation by the
gyroscope chain and refits translation to the same corners, so that a network
trained on these labels receives the sensor that helped create its target. We
compare its held-out $R^2$ with the same architecture trained on unmodified,
cleaned PnP labels. Gyroscope-constrained labels are not used for any
reported model.
